\documentclass[10pt,a4paper]{article}
\usepackage[T1]{fontenc}
\usepackage[utf8]{inputenc}
\usepackage{lmodern}
\usepackage{amsmath,amssymb}
\usepackage[margin=21mm]{geometry}
\usepackage{microtype}
\usepackage[hidelinks]{hyperref}
\usepackage{enumitem}
\setlist[itemize]{leftmargin=*,nosep,topsep=3pt}
\setlength{\parindent}{0pt}
\setlength{\parskip}{5pt}
\setlength{\emergencystretch}{3em}
\newcommand{\question}[1]{\par\medskip{\normalsize\bfseries #1}\par\smallskip}
\pagestyle{plain}

\begin{document}
{\Large\bfseries Conditions for the Generation of Beta Oscillations in the Subthalamic Nucleus-Globus Pallidus Network\par}

\textbf{Paper:} Nevado Holgado, Terry and Bogacz (2010).\quad Delayed population-rate modelling study.\par
\href{https://doi.org/10.1523/JNEUROSCI.0817-10.2010}{DOI: \nolinkurl{10.1523/JNEUROSCI.0817-10.2010}}\par
{\footnotesize Page references below use PDF page numbers, including front matter.}

\question{To which general question does this paper relate?}
A concrete delayed nonlinear model and conditions for oscillation generation. Relevant to the notebook's model class; its mechanistic predictions require separate behavioural validation.

\question{How does it answer the general questions?}
\textbf{Model mapping.} With rates $r_S,r_G$ and constant cortical/striatal drives,
\[
\begin{aligned}
\tau_S\dot r_S&=-r_S+F_S(-w_{GS}r_G(t-T_{GS})+w_{CS}\mathrm{Ctx}),\\
\tau_G\dot r_G&=-r_G+F_G(w_{SG}r_S(t-T_{SG})-w_{GG}r_G(t-T_{GG})-w_{XG}\mathrm{Str}).
\end{aligned}
\]
Take $x=r-r^*$, $a_i=-1/\tau_i$, $b^x=0$, and signed nonlinear input rows $(0,-w_{GS})$, $(w_{SG},-w_{GG})$. Set $f_i(v)=[F_i(q_i^*+v)-F_i(q_i^*)]/\tau_i$. The reported sigmoid has $0\leq F_i'\leq1$, giving global slope and sector bounds $[0,1/\tau_i]$ after centring (derived from their Eq.~3).

\textbf{Parameters.} $\tau_S=6$, $\tau_G=14$, $T_{SG}=T_{GS}=6$, $T_{GG}=4$ in ms. With $F_i(q)=M_i/[1+e^{-4q/M_i}(M_i-B_i)/B_i]$, use $(M_S,B_S,M_G,B_G)=(300,17,400,75)$ and $(\mathrm{Ctx},\mathrm{Str})=(27,2)$ in spikes/s. The weight vector $(w_{SG},w_{GS},w_{GG},w_{CS},w_{XG})$ is $(19,1.12,6.60,2.42,15.1)$ for the healthy fit and $(20,10.7,12.3,9.2,139.4)$ for the diseased fit (Tables~1, 3; PDF pp.~2--4). Their pulse experiments use a fitted additive electrode-current term in the STN equation (Eq.~5); no LFP observation map is identified.

\textbf{Interpretation.} Delays, coupling and external drive determine oscillation onset within this model. Predictions about removing beta are model-level results, not proof that beta suppression alone treats symptoms (PDF pp.~5--12).

\question{What papers does it cite that could provide more information?}
\begin{itemize}
\item Plenz and Kitai (1999), \emph{A basal ganglia pacemaker formed by the subthalamic nucleus and external globus pallidus}: experimental loop hypothesis (PDF p.~13).
\item Terman et al. (2002), \emph{Activity patterns in a model for the subthalamopallidal network of the basal ganglia}: alternative spiking mechanism (PDF p.~13).
\end{itemize}

\end{document}
